\chapter{Bad Counterfoils: A General Theory of Borrowed Authority}
\label{ch:bad-counterfoils}

\begin{plainlanguage}
\textbf{In plain terms:} One idea from Chapter~5 — that mistakes happen when low-quality evidence gets treated as high-quality — turns out to explain failures everywhere: institutions that declare success without changing, moral panics, exaggerated warnings. This chapter states that pattern directly as the book's single most exportable idea, useful even to a reader who rejects everything else here.
\end{plainlanguage}

\begin{chapternote}
A note on sequence. This part was drafted with full knowledge of Part~V's actual outcome (Chapters~14--15: a partial, not clean, transport of the recursive compression conjecture to creativity, and an explicitly unfinished compression-efficiency conjecture). That outcome is referenced directly below where relevant. The point made here is not that Part~V's result was unknown --- it is that this part's central argument does not depend on which way Part~V came out. The argument was already complete by the end of Part~II, and Part~V's partial, imperfect result is, if anything, a better test of that independence than a clean success would have been: the reader can check, concretely, that Chapter~17's conclusion does not lean on Part~V having gone better than it did.
\end{chapternote}

\section{One sentence doing more work than it looked like}

Chapter~\ref{ch:fidelity} produced a sentence that was supposed to be a summary, not a discovery in its own right: intelligence-as-self-description fails when low-fidelity counterfoils are mistaken for hard eliminations. Across the chapters since, that sentence has kept reappearing under different names. Chapter~\ref{ch:room-for-river} found institutions declaring closure ($N(t)=1$) without having operationally achieved it ($O(t)=0$) --- a nominal elimination mistaken for a real one. Chapter~\ref{ch:development} found a developmental literature in which production-level overextension looked, from the outside, like a recognition failure, when comprehension underneath it was already intact --- an apparent boundary mistaken for a real one. Chapter~\ref{ch:counterfoil} found an entire cultural pedagogy organized around manufacturing low-stakes, high-margin counterfoils precisely so a learner could practice telling forced eliminations apart from real ones before the stakes got serious. The pattern was never confined to intelligence. This chapter states it on its own terms.

\section{The general theory}

\begin{definition}[Borrowed Authority]
Let $q_i$ be the source-independence component of fidelity introduced in Chapter~\ref{ch:fidelity} --- the degree to which an elimination's authority survives being stripped of the social standing of whoever asserts it. Define
\[
B_i = 1 - q_i,
\]
the \emph{borrowed-authority score} of constraint $i$: how much of its apparent force depends on the teller rather than the world.
\end{definition}

\begin{definition}[Counterfoil Distortion Index]
Let $D_i$ be the ratio of an eliminating claim's depicted severity to its actual severity:
\[
D_i = \frac{\text{depicted severity}}{\text{actual severity}}.
\]
A large $D_i$ indicates rhetorical amplification --- the claim is dramatizing an outcome beyond what the evidence supports, which corresponds to low severity calibration ($s_i$, from Chapter~\ref{ch:fidelity}) for the same constraint.
\end{definition}

\begin{proposition}[checked for circularity]
World-forced eliminations are characterized by high overall fidelity $w_i$; rhetorical eliminations are characterized specifically by high $B_i$.
\end{proposition}

Stated this baldly, the proposition risks being empty rather than informative, and it is worth being explicit about exactly where the emptiness would come from rather than letting a reader discover it unaided. Since $q_i$ is one of the four multiplicative factors composing $w_i$ ($w_i = d_i \cdot s_i \cdot c_i \cdot q_i$, in the notation fixed in Chapter~\ref{ch:fidelity}), an elimination with low $q_i$ will tend to have low $w_i$ as well, simply by the arithmetic of the formula --- unless the other three factors happen to compensate. ``Rhetorical eliminations have high $B$ and low $w$'' is, read this way, close to a restatement of how $B$ and $w$ are defined in terms of the same underlying $q$, not an independently derived empirical finding.

The chapter's actual, non-circular claim is narrower and more interesting than the proposition as first stated, and it is the claim worth defending: in the cases examined across this book --- the empty-bowl pedagogy of Chapter~\ref{ch:counterfoil}, the rhetorically amplified institutional myths this chapter goes on to discuss --- low $q_i$ tends to co-occur with low $d_i$, low $s_i$, and low $c_i$ \emph{simultaneously}, rather than independently. The empty bowl is not merely authority-dependent (low $q$); it is also low-directness (poverty does not follow deterministically from a single act of childhood laziness), low-severity-calibration (the depicted certainty of ruin vastly overstates the actual conditional probability), and low-closure (real recovery paths --- charity, family support, second chances --- exist and are simply not shown). This joint co-occurrence across all four components is what the chapter claims, and it is an empirical pattern about how low-fidelity counterfoils are actually constructed in practice, not a logical consequence of how the four components happen to be defined. It is also, importantly, not established with statistical rigor here --- it rests on the worked examples examined in Chapter~\ref{ch:counterfoil} and the case sketched below, and should be read with the same caution Chapter~\ref{ch:two-domains} applied to its own two-domain recurrence: suggestive, not demonstrated.

\section{Exporting the machinery}

Political ideologies, religious systems, moral panics, risk communication, and institutional self-justification all run on collections of counterfoils --- exaggerated, rhetorically amplified non-options offered to make a predetermined conclusion feel self-evident, structurally indistinguishable in form from the empty bowl, differing only in the domain and the stakes.

Take a single brief case to show the machinery travels rather than merely asserting that it does. A moral panic --- the recurring pattern in which a specific, often statistically rare harm is presented as an imminent, near-certain threat to children or to social order, justifying urgent restriction of some named activity --- can be scored against all four fidelity components directly. Causal directness is typically low: the panic's narrative usually requires several unstated, contingent intervening steps between the named activity and the depicted harm. Severity calibration is typically low: $D_i$, the ratio of depicted to actual severity, tends to run far above $1$, since panics are, definitionally, panics partly because the depicted risk outpaces the measured one. Closure is typically low: panics rarely acknowledge the recovery paths, safeguards, or base-rate context that would weaken their force. And source independence is typically low: panic-narratives draw heavily on the authority of the messenger (an institution, a movement, a media figure) rather than on claims that would retain their force if asserted by an unaffiliated stranger. The same four-component audit that distinguished the stove warning from the empty bowl in Chapter~\ref{ch:counterfoil} distinguishes a documented public-health risk communication from a moral panic dressed in its vocabulary --- and does so by the same mechanism, applied outside intelligence entirely.

\subsection*{A second, more contested case: deflection}

A second case is worth including specifically because it is harder than the first, and the difficulty is instructive about a limit of this chapter's machinery rather than a flaw in it. At the same ``Is AI a Threat?'' panel already cited in Chapters~\ref{ch:room-for-river} and \ref{ch:two-axes}, panelist Jens Munch, a technology and AI strategist, entrepreneur, and investor, argued that public and policy attention to speculative risks from rogue superintelligence or artificial general intelligence functions as a convenient deflection from a more mundane, more immediately tractable problem: as AI systems are continuously updated and re-versioned, legal liability for their decisions becomes structurally difficult to assign, eroding corporate accountability and democratic oversight in ways current legal frameworks are not equipped to track.

The structure of this claim is exactly the structure Chapter~\ref{ch:counterfoil} introduced and this chapter has been auditing: a vivid, high-salience narrative (catastrophic superintelligent risk) is alleged to be drawing attention and resources away from a less dramatic but more directly actionable structural problem (legal-accountability erosion), in a manner the four fidelity components could in principle assess --- is the attention-capturing narrative's severity calibrated against its actual probability, does it depend on contingent and currently unrealized technical developments rather than direct, present mechanisms, and how much of its force depends on the prominence of the people discussing it rather than on demonstrated likelihood?

This book takes no position on how that audit would come out, and the case is included specifically to mark where the machinery's confident application from the previous case stops being available. Unlike the moral panic example, whether AGI risk is itself overstated, understated, or roughly proportionate to the attention it receives is a genuinely contested empirical and technical question among serious, informed people, not a case where one side's low fidelity is already established by available evidence. What can be said without resolving that question is narrower and still useful: the deflection structure being alleged --- one salient risk narrative crowding out attention to a less salient but more tractable one --- is a real and recurring pattern in how institutions and publics allocate limited attention, independent of whether superintelligence risk specifically is the salient narrative doing the crowding-out in this instance. The fidelity audit is the right tool for evaluating that structural claim case by case; it does not, by itself, settle the case in advance, and a book whose central caution is against treating borrowed authority as a substitute for demonstrated mechanism should not exempt its own use of a single contested example from that same caution.

This is offered as the book's most exportable single idea: a reader who rejects everything else argued here --- the recursive compression principle, the weak-realist resolution of Chapter~\ref{ch:resolving-fork}, even the relevance of intelligence as the book's organizing example --- could still take the four-component fidelity audit and use it on a piece of rhetoric encountered tomorrow. Good reasoning, in any domain, consists substantially in learning which eliminations deserve authority. Bad reasoning consists, with notable regularity, in mistaking a constraint that is low on directness, low on calibration, low on closure, and high on borrowed authority for one that is none of those things.
